% Gauss, Werke X/1 (1917), native108–112, printed101–105.
\SourceLeaf{108}\ArithmeticHead{101}{R}{PFAFF AN GAUSS.}
\NumberedItem{2}
\begin{center}\large\textsc{Pfaff} an \textsc{Gauss}.\end{center}
\noindent\makebox[\linewidth][s]{. . . . . . . . . . . . . . . . . . . . . . . . . . .}\par
\noindent\hspace*{\parindent}\{Es war mir erfreulich, dass Sie meine wenigen unbedeutenden Anmerkungen
gut aufgenommen haben, und selbst noch mehrere wünschen. Es gehört zu den Vorzügen der
Mathematik, dass darin nicht viel zu disputiren ist, und dass man von dem unangenehmen Eindruck,
welchen die Rechthaberey erweckt, selten etwas erfährt. — Ihr Verlangen nach mehreren Bemerkungen,
und nach Angabe der Stellen, die etwa deutlicher seyn könnten, kann ich übrigens jetzt nicht so
wie ich wünschte erfüllen, da ich dazu Ihre Abhandlung, die mir inzwischen aus dem Gesicht
gekommen war, wieder genauer durchgehen müsste, als es meine Zeit zulässt, indem ich 3 Collegia
zu lesen, u. ausserdem noch eigne Beschäftigungen habe, welche ich um so weniger jetzt aussezen
kan, da ich auf Michaelis eine Reise in mein Vaterland zu machen gedenke. —

Einige Kleinigkeiten, die mir aufgestossen waren, wovon ich das vorigemal etwas erwähnte
(meistens Schreib- und Sprachfehler) halte ich jetzt nicht einmal für nöthig anzuzeigen, theils
weil Ihre Bemerkung wegen der sin. u. cos., dass die Leser sich selbst daraus finden werden,
auch hier gälte, theils weil Sie ohnehin das Ganze vor dem Abdruck noch einmal selbst genau
durchsehen werden. — Wenn ich äusserte, dass Ihre Abhandlung noch hie u. da an Deutlichkeit
gewinnen könnte, so folgte ich dabey dem Total-Eindruck, den sie bey dem ersten Lesen auf mich
gemacht hatte, u. vermuthlich auch auf andere Leser machen wird: allerdings liegt aber viel
davon an der Kürze, die Sie, wie Sie sagen, absichtlich gewählt haben. Dass Sie für die letzte
Figur die Gleichung und damit ein concretes Beyspiel geben wollen, wird, wie ich auch glaube,
zur Erläuterung dienen: ich dachte erst, dass es noch dienlicher wäre, wenn Sie, nicht am Ende
in einer Anmerkung, sondern schon vor der allgemeinen Untersuchung ein solches Beyspiel
entwickelten, zur Vorbereitung auf das folgende; doch mag Ihnen diess für Ihren Zweck zu
weitläufig scheinen. Ich hatte mir jenes etwa zwischen § 16 u. § 18 gedacht, an der Stelle
von § 17. Sie sagen in Ihrem Brief, dass Sie den ganzen \underline{§ 18} bloss zur leichteren
Verständlichkeit vorläufig beygefügt haben, u. anfangs
\par\clearpage
\SourceLeaf{109}\ArithmeticHead{102}{L}{ALGEBRA.\quad BRIEFWECHSEL.}
\noindent willens gewesen seyen, ihn ganz wegzulassen. Ohne Zweifel meynen Sie damit den
\underline{§ 17}, u. ich muss gestehen, dass mir dieser § nicht recht hat einleuchten wollen,
besonders, da man fragen könnte, warum die niedrigen Potenzen von $r=\infty$ gegen die ersten
weggelassen werden, da dieser ihr Coëfficient $=0$ wird (welcher Zweifel übrigens durch
Betrachtungen der Gleichungen zwischen $y$ u. $x$ gehoben werden möchte). Was den § 18 betrifft,
so habe ich bey wiederholter Ansicht leicht bemerken können, dass die Puncte (1) (3) (5)
.\,.\,. nicht in der Figur 2 abgebildet sind (welches schon die Menge der Puncte in der Figur,
nur bis 15, hätten zeigen können). Indessen war doch diese Auslegung, welche sonst dem übrigen
Sinn gar keinen Eintrag thut, ziemlich natürlich, da die Worte § 18 »Designato puncto
circumferentiae .\,.\,. per (1)« auf die Figur hinleiten, u. man eher geneigt ist, unter geraden
Linien sich Schenkel von Winkeln, als ohne Erinnerung krummlichte Schenkel dadurch angedeutet
zu denken. Wenn Sie, wie Sie schreiben, zur Verhütung alles Missverständnisses noch etwas
beyfügen, so wird es desto besser seyn. — Auch die Anmerkg. zu § 20 ist in dieser Hinsicht
etwas dunkel, und wohl auf eine oder die andere Art verschrieben. Bey eben diesem § 18 steht
als Note »Suppono summam $S$ \underline{majorem} esse quam $\surd\frac12$; alioquin radium $R$
\underline{majorem quam} 1 accipere oportet«. Es wird hier etwas verschrieben seyn, u. ohne
Zweifel heissen müssen: $S$ \underline{minorem} quam $\surd\frac12$. Dabey würde aber die
Bedenklichkeit entstehen, ob u. warum diese einschränkende Bedeutung zulässig sey. Dass die
Linie, welche $=1$ gesetzt wird, willkührlich ist, hat natürlich hierauf keinen Einfluss.
Ferner sagen Sie: summa
\[
A\sin(m-1)\varphi+\frac BR\sin(m-2)\varphi+\frac C{R^2}\sin(m-3)\varphi
+\frac D{R^3}\sin(m-4)\varphi+\mathrm{ec.}
\]
\noindent certo non potest esse major quam $S\,(=A+B+C+D+\mathrm{ec.})$.
Dabey kann man wieder nach dem Grund fragen, da Sie $R<1$ nehmen, also
$\frac1R,\ \frac1{R^2},\ \frac1{R^3}$ ec. grösser als 1 sind.
Übrigens weiss ich nicht, warum $R$ nicht auch grösser als 1 sollte genommen werden dürfen
(der sin. tot. ist zwar $=1$, kommt aber natürlich hier nicht in Betrachtung).
Für $R>1$ würden dann beyde Bedenklichkeiten wegfallen. Aus diesen Gründen, so weit ich die
Sache für jetzt übersehe, scheint mir daher eine kleine Änderung nöthig zu seyn, die sich leicht
anbringen liesse. —
\par\clearpage
\SourceLeaf{110}\ArithmeticHead{103}{R}{PFAFF AN GAUSS.}
\noindent\hspace*{\parindent}Ihre Darstellung des wirklich sinnreichen u. einfachen Beweises von
\textsc{d’Alembert} ist fast etwas zu gedrängt ausgefallen, so dass es dem Leser, der die
Abhandlung nicht selbst hat, schwer seyn möchte, sich sogleich darein zu finden. Vielleicht
liessen sich die Hauptpuncte, worauf es dabey ankommt (nehmlich der Ausdruck von $x$ durch
eine convergirende Reihe nach $X$, für kleine Werthe von $X$, u. der Schluss von kleinen
Werthen auf grössere) etwas fasslicher herausheben. Inzwischen ist mir noch eine Bedenklichkeit
bey \textsc{d’Alemberts} Beweis aufgestossen. Wenn nehmlich nach \textsc{Bougainville} eine
allgemeine algebraische Gleichung zwischen $x$ u. $X$ angenommen wird, so lässt sich zwar
(mit Hülfe des \textsc{Newton}schen \#ogramms) $x$ durch eine Reihe nach Potenzen von $X$
darstellen, aber für die Coëfficienten dieser Reihe ergeben sich öfters selbst
\textls[120]{höhere Gleichungen}, daher jene auch \textls[120]{unmöglich} ausfallen können,
wo dann nicht ohne Beweis vorausgesetzt werden dürfte, dass sie von der Form $a+b\surd-1$
sind. So möchte der Beweis in seiner \textls[120]{Allgemeinheit sogleich} vorgetragen eine Art
von logischem Zirkel enthalten. Doch liesse sich auf diese Einwendung vielleicht noch antworten.
— Eben indem ich diesen Brief schreibe, bin ich, da ich mir \textsc{d’Alemberts} Beweis auf eine
andre Art entwickeln wollte, noch auf eine Schwierigkeit gestossen, ich habe aber jetzt nicht Zeit,
der Sache weiter nachzudenken. Doch lässt sich vielleicht schon aus folgendem einfachem Beyspiel
übersehen, worauf es ankommt, u. was ich damit meyne. Es sey nehmlich die quadratische Gleichung:
$y=Ax-Bx^2$, so wird $x$ durch eine \textls[120]{convergirende} Reihe nach $y$ dargestellt
(u. zugleich möglich) wenn $\frac yB<\frac{A^2}{4B^2}$. Nun könnte man neue Werthe von $y$ u. $x$
annehmen, nehmlich $y'=y+\eta$, $x'=x+\chi$, u. zwey Gleichungen formiren, die eine zwischen $y$
u. $x$, die andre zwischen $\eta$ u. $\chi$. Da nun für die erste, $y$ die vorerwähnte Grenze
hat, u. für die andre $\eta$ auch eine gewisse Grenze, so könnte man schliessen, $y+\eta=y'$
werde eine \textls[120]{grössere} Grenze haben, als die vorige, u. so immer weiter.
Dieser Schluss wäre aber natürlich unrichtig, wie sich auch ergibt, wenn man die Gleichung
für $\eta$ entwickelt.

Sie ist:
\[
\eta=(A-2Bx)\chi-B\chi^2
\]
\noindent also ist die Grenze für sie:
\[
\frac\eta B<\frac{(A-2Bx)^2}{4B^2}.
\]
\clearpage
\SourceLeaf{111}\ArithmeticHead{104}{L}{ALGEBRA.\quad BRIEFWECHSEL.}
\noindent Aber aus der ersten Gleichung ist
\[
x=\frac A{2B}\pm\surd\left(\frac{A^2}{4B^2}-\frac yB\right)
\]
\noindent also
\[
\frac{(A-2Bx)^2}{4B^2}=\frac{A^2}{4B^2}-\frac yB,
\]
\noindent mithin
\[
\frac\eta B<\frac{A^2}{4B^2}-\frac yB,
\]
\noindent woraus weiter nichts folgt als
\[
\frac\eta B+\frac yB\ \text{d. i.}\ \frac{y'}B<\frac{A^2}{4B^2},
\]
\noindent die vorige Grenze. Überhaupt ist die Kleinheit der Werthe von $x$, für welche
die Reihe convergirt, nicht \textls[120]{absolut}, sondern hängt von den Coëfficienten ab,
und diese \textls[120]{ändern sich mit der Gleichung}. Die allgemeinere Entwicklung dieser
Bemerkungen u. ihres näheren Zusammenhangs mit \textsc{d’Alemberts} Schlüssen, worauf sie sich
beziehen, kann ich jetzt nicht übernehmen: inzwischen würde es mich freuen, wenn Sie solche
(so wie die vorige Einwendung gegen \textsc{D’Al.} Beweis) einer weiteren Überlegung werth hielten. —

Ich hatte den Band der Berliner Memoires von 1772 gerade im Haus, aber die Abhandlung von
\textsc{La Grange} über die \textls[120]{unmögl. Wurzeln} noch nicht gelesen, an dem Post-Tag,
da mir gemeldet wurde, dass Sie solche zu sehen wünschten. Wegen des für Sie wichtigen
arithmetischen Inhalts wollte ich Sie nicht warten lassen, u. las also gleich den Aufsatz von
\textsc{La Grange} so weit durch, als nöthig war, um die Haupt-Momente deutlich zu übersehen,
besonders da die Sache mit vieler Klarheit vorgetragen ist. Sie werden aus dieser Abhandlung
sehen, dass \textsc{La Grange} die Lücken in \textsc{Eulers} Beweis auszufüllen sucht, besonders
was theils den von \textsc{Foncenex} als nicht evident bemerkten Saz, theils die Bestimmung
möglicher Werthe von $\alpha,\ \beta$ ec. durch $u$ (§ 7 Ihrer Dissert.) betrifft.
In Rücksicht des 2\textsuperscript{ten} Punctes möchte es noch darauf ankommen, ob die Säze
von der Elimination in der früheren Abhandlung von \textsc{La Grange} (die bekanntlich auch
von \textsc{Michelsen} deutsch übersezt ist) völlig streng erwiesen sind. Besonders ist mir
dabey folgende Bedenklichkeit eingefallen: Wenn $r$ Gleichungen für $r$ unbekannte Grössen
$u,\ \alpha,\ \beta,\,.\,.\,.\ \lambda,\ \mu\,.\,.\,.$ gegeben sind, so lassen sich
\textls[120]{im allgemeinen} immer $\alpha,\ \beta,\,.\,.\,.\ \lambda,\ \mu,\,.\,.\,.$
\par\clearpage
\SourceLeaf{112}\ArithmeticHead{105}{R}{PFAFF AN GAUSS.}
\noindent rationell durch $u$ ausdrücken (wie unmittelbar aus dem Verfahren der Elimination
folgt), aber es können die Werthe dieser Grössen unbestimmt $\alpha=\frac00$ werden, nicht
nur für \textls[120]{gleiche} Werthe aus der Gleichung für $u$, sondern für \textls[120]{jedes}
$u$ (wo man dann eigentlich keine Gleichung des ersten Grads für $\alpha$ erhält).
Doch möchte eine nähere Betrachtung der Elimination für den Fall, der hier zunächst in
Betracht kommt, diese Bedenklichkeit heben.

Ich habe Sie schon zu lang mit diesem Brief aufgehalten. Ich will also nichts weiter beyfügen,
als dass mir fernere Nachrichten und Mittheilungen von Ihnen jederzeit sehr willkommen seyn
werden, u. dass ich mit der aufrichtigsten Hochachtung verharre
\par\medskip
\noindent\begin{minipage}[b]{.35\linewidth}
\begin{center}Helmstedt,\\8. Jul. 1799\end{center}
\end{minipage}\hfill
\begin{minipage}[b]{.5\linewidth}
\begin{center}Ihr\\ergebenster Freund\\J. F. \textsc{Pfaff}.\end{center}
\end{minipage}\par\medskip
\noindent\hspace*{\parindent}N.S. Sollte in Braunschweig oder Wolfenbüttel \textsc{Cheyne}’s
\textit{Methodus fluxionum}[*)] zu haben seyn, so würden Sie mich durch dessen Mittheilung
auf kurze Zeit verbinden. H. H[of]r[ath] \textsc{Zimmermann} bitte ich mich geh[orsam]st zu
empfehlen, wie auch H. Bergrath \textsc{Volckmar}, dessen nähere hier kürzlich gemachte
Bekanntschaft mir sehr angenehm gewesen ist. Noch fällt mir eben jetzt folgendes bey,
das vielleicht doch auch eine Erwägung verdiente. Die Gleichungen $T=0$ u. $U=0$ nach $x$ u. $y$
entwickelt gäben eliminiert Gleichungen für $x$ u. für $y$, deren jede eine mögliche Wurzel haben
müsste. Sollte diess nicht aus bekannten Säzen von den möglichen Wurzeln bey Gleichungen vom
$2r^{\mathrm{ten}}$ u. $2r-1^{\mathrm{ten}}$ Grad herzuleiten seyn? Indessen möchte doch der
Beweis vielleicht nicht leichter seyn, als bey \textsc{Foncenex}’s Verfahren.
Übrigens wäre eine nähere Betrachtg. der Gleichungen für $x$ u. $y$ immer gut.
Bey \textsc{Foncenex} muss $u$ nicht möglich seyn, aber hier \textls[120]{müssen} $x$ u. $y$
möglich seyn. Doch ich will Sie nicht länger mit flüchtigen Gedanken aufhalten.\}
\par\smallskip\noindent\rule{.25\linewidth}{.4pt}\par\smallskip
\begingroup\small
\noindent\hspace*{\parindent}[*) \textsc{George Cheyne}, \textit{Fluxionum methodus inversa sive
quantitatum fluentium leges generaliores,} Lond. 1704.]
\par\endgroup
\vfill\noindent X\textsubscript{1}.\hfill 14
\clearpage
